% !TeX root = ../../protocolo.tex

In this chapter the advances on the substrate effects in the optical properties of metasurfaces are presented. This results consist in the comparison between experimental results  measured at the Instituto Nacional de Astrofísica, Óptica y Fotónica (INAOE), where our collaborators have the capacity to synthesize metasurfaces of spherical gold nanoparticles (AuNP) on a Corning glass substrate by dewetting a gold continuous film. In addition, the experimental  setups presented in Section \ref{s:Setup} are mounted at INAOE, where the measurements of extinction and  reflectivity spectra  were performed for two samples of metasurfaces. The obtained experimental spectra are discussed in the following  section with an analysis considering the Coherent Scattering Model discussed in Section \ref{s:CSM} and Finite Elment Method (FEM) simulations performed in  COMSOL Multiphysics$^\text{TM}$ version 6.1, from which each contribution of the substrate effects  on the optical response of the metasurfaces are estimated. 

The metasurfaces analyzed in this chapter were synthesized from films with  an initial thicknesses $d = (3.0\pm 0.1)$~nm and $d = (5.0\pm 0.2)$~nm which were held inside a Caisa DTT434 oven at $550^\circ$C for 3 hours.  Images of Scanning  Electron Microscopy (SEM) of each  metasurface are shown in the background of Figs. \sfref{fig:morpho}{a)} and \sfref{fig:morpho}{b)} for the $d=3$~nm and $d=5$~nm samples, respectively. The orange curves in the foreground correspond to the histograms of the AuNPs classified by their equivalent radius, while the red and green curves correspond to the Gaussian-fitted distributions of the histograms, from which the mean radius $\mu$ and the standard deviation $\sigma$ of the size distribution of each sample were determined. The sample with initial thickness of $d=3$ nm, shown in Fig. \sfref{fig:morpho}{a)}, is characterized with parameters $\mu = (5.78 \pm 0.06)$~nm, $\sigma = (1.89\pm 0.06)$~nm, and a surface coverage $\Theta = 0.20~\pm~0.01$. In contrast, the sample with initial thickness of $d=5$ nm, shown in Fig. \sfref{fig:morpho}{b)}, is characterized with $\mu = (7.30 \pm 0.05)$~nm, $\sigma = (2.30\pm 0.05)$~nm, and $\Theta = 0.22~\pm~0.01$. 

\begin{figure}[h]\centering
	\includegraphics[scale = 1]{non-int/Fig-Sample.pdf}
	\caption{%
		SEM image and histogram of the equivalent radius of the AuNPs in a metasurface fabricated from a gold film with an initial thickness \textbf{a)}  $d = (3.00\pm 0.02)$~nm and  \textbf{b)} $d = (5.00\pm 0.02)$~nm, both showing a normal distribution of the equivalent radius of the AuNPs with a mean $\mu$, standard deviation $\sigma$, and surface coverage $\Theta$.}
	\label{fig:morpho}
\end{figure}

